\section{Related Work}
\label{sec:related}

%Before we lay out our approach, it helps to see exactly what has already been tried on this problem, and where each attempt stops short.

\textit{Defects4J} \citep{just2014defects4j} is one of the most widely used benchmark datasets for software debugging and defect analysis research. It provides a curated collection of reproducible real-world bugs from 17 open-source Java projects, along with their original bug reports, reproducible failures, buggy and fixed source code, and buggy-to-fixed patches. This rich and reproducible set of artifacts makes \textit{Defects4J} particularly suitable for automated defect classification. In this research, we leverage these artifacts to investigate the classification of \textit{Defects4J} bugs into ODC defect categories using large language models (LLMs). The remainder of this section provides the necessary background on \textit{Defects4J}, ODC, and LLM-based defect classification.

\subsection{Existing \textit{Defects4J} Bug Type Classification}

%Prior work on classifying \textit{Defects4J} bugs falls into three lines, each covering a different slice of the dataset with a different goal.

Van der Spuy et al. \citep{vanderspuy2025anatomy} manually classified 488 faults using control-flow and data-flow graph abstractions, identifying six control-flow and two data-flow categories. This structural perspective provides insight into how defects alter program execution; however, it does not necessarily capture the underlying reason for the defect. For example, a missing null check and an omitted feature may result in similar control-flow changes while representing fundamentally different defect causes.

Sobreira \textit{et al.} \citep{sobreira2018dissection} analyzed 395 bugs from \textit{Defects4J} v1.2, focusing on repair actions and patch patterns. Their findings showed that a substantial proportion of patches involve modifications to conditional blocks. Although patch-based analysis is valuable for understanding how developers repair defects, patch characteristics primarily describe the changes made to the program rather than the underlying defect mechanism that necessitated those changes.

Xuan \textit{et al.} \citep{xuan2017better} manually examined 50 randomly selected bugs to investigate test-suite quality and fault fixability. While their analysis provides detailed insights into individual defects, the relatively small sample limits the generalizability of the findings to the broader \textit{Defects4J} dataset.

These studies provide valuable perspectives on defect structure, repair actions, and fixability. However, they examine relatively limited subsets of \textit{Defects4J}, rely primarily on code-structural or patch-based taxonomies, and do not apply a semantic defect-type taxonomy to the full set of bugs. Our work addresses this gap by investigating ODC-based semantic defect classification across the \textit{Defects4J} dataset using large language models.

\subsection{Orthogonal Defect Classification (ODC)} \label{subsec:ODC}

ODC was introduced by Chillarege \textit{et al.}\ \citep{chillarege1992odc} as
a framework that classifies defects along orthogonal (Independent) attributes. It was
originally designed to give an in-process signal during active development. The ODC framework defines eight attributes organized into two phases: \emph{Opener} attributes at discovery time and
\emph{Closer} attributes at fix time, summarized in Table~\ref{tab:odc-attributes}.

\begin{table}[H]
\centering
\caption{ODC Opener and Closer attributes \citep{ibm2013odc}.}
\label{tab:odc-attributes}
\setlength{\tabcolsep}{5pt}
\begin{tabularx}{\columnwidth}{@{}llX@{}}
\toprule
\thead{Phase} & \thead{Attribute} & \thead{Description} \\
\midrule
\multirow{3}{*}{Opener} & Activity    & Specific activity being performed when the defect was found (e.g.\ code inspection or unit test) \\
                         & Trigger     & Condition that exposed the defect \\
                         & Impact      & User-visible consequence of the defect \\
\midrule
\multirow{5}{*}{Closer} & Defect Type & The actual correction made \\
                         & Target      & Design or Code \\
                         & Qualifier   & Missing, Incorrect, or Extraneous \\
                         & Age         & Whether the faulty code is base, new, rewritten, or refixed \\
                         & Source      & Whether the faulty code is developed in-house, reused from a library, outsourced, or ported \\
\bottomrule
\end{tabularx}
\end{table}

The central attribute for our work is the Defect Type, which represents the semantic nature of the actual code correction made to fix a bug and define seven defect type categories  detailed in Table~\ref{tab:}.

\textcolor{red}{Table representing the following info: 1st column - Control and Data Flow (the first four) and Structural (the last three) 2nd column:  Algorithm/Method, As
signment/Initialization, Checking,Timing/Serialization,Function/Class/Object,Interface/O
OMessages,andRelationship. 3rd column : Description}


ODC defect types are actively used in modern software debugging research. Pan e al.\citep{esfahani2024understanding} applied ODC to classify defects in LLM-generated code in 2024, demonstrating that the taxonomy generalizes beyond human-written software. The official ODC v5.2 specification from IBM \citep{ibm2013odc} continues to serve as the reference standard for the framework.

%\subsection{\textcolor{red}{mapping}}
%Going through the eight attributes from Table~\ref{tab:odc-attributes} shows how much of each one Defects4J can actually support.

%\begin{itemize}
%\item \textbf{Activity.} Defects4J bugs are reproduced by running a unit test, so Unit Test is the natural fit for this attribute. Bug reports sometimes describe how the problem was noticed, which could point to a different activity such as inspection, but this shows up inconsistently across reports.
%\item \textbf{Trigger.} The same reasoning applies here. The unit test that reproduces the bug is the closest available signal for the condition that exposed it, and a report's own description could suggest something else, again inconsistently.
%\item \textbf{Impact.} This is what we derive directly from the bug report and the failure behavior, one of the two attributes this study actually classifies.
%\item \textbf{Defect Type.} This is what the study tries to predict from the buggy code, then checks that prediction once the fixed code becomes available.
%\item \textbf{Target.} By definition this is Code. Defects4J bugs are fixed in source, and there is no separate design document ODC's Design target would require.
%\item \textbf{Qualifier.} This becomes available once the fix itself is visible, since it describes whether the fix supplied something missing, corrected something wrong, or removed something unneeded.
%\item \textbf{Age.} Defects4J bugs come from real project histories with full commit logs, so age could plausibly be traced through the commit history of the affected code. This is not something pursued in this study.
%\item \textbf{Source.} Most Defects4J projects are long running, community maintained Java libraries, so Source would intuitively be Developed In-House for nearly all of them. This is a reasonable expectation, though unconfirmed.
%\end{itemize}

%That leaves Defect Type and Impact as the two attributes this study classifies. The others fall into three groups: fixed by how the dataset is built, a reasonable guess without confirmation, or visible only once the fix itself is known.

\subsection{LLM-based Defect Classification}

Recent studies have demonstrated that Large Language Models (LLMs) can effectively classify software defects when provided with relevant contextual information.

\citep{colavito2024llm} showed that LLMs can classify issue reports (closely related to bug reports) with reasonable accuracy, validating the basic hypothesis that LLMs are suitable for defect classification tasks. \citep{koyuncu2025exploring} extended this by using LLMs and several prompt engineering techniques, including few-shot, chain-of-thought, tree of thoughts, and plan-and-solve prompting, to achieve fine-grained bug report categorization. 

\citep{li2024knowbug, du2024llm} uses LLMs, leveraging several prompting strategies to classify bug reports into four categories based on their behavior and reproducibility. These studies showed that large language models benefit from structured prompting strategies, ranging from simple task prompts to reasoning‑based approaches like Chain‑of‑Thought and Tree‑of‑Thought, enriched with few‑shot examples and domain knowledge integration, to improve accuracy, generalization, and interpretability of bug report categorization.

Relatedly, Kang \textit{et al.} \citep{kang2023autosd} incorporated Zeller's scientific debugging method \citep{zeller2009why} into LLM prompting through AutoSD. In their approach, the model formulates a hypothesis, performs a debugging experiment to test it, and uses the resulting evidence to determine whether the hypothesis is supported, primarily for explainable fault localization. Building on this idea, our iterative reasoning process (hypothesis--predict--experiment) uses targeted probes of the available evidence to evaluate competing hypotheses and ultimately classify the defect according to ODC type.

Compared with traditional machine learning approaches, LLMs offer a key advantage in their ability to jointly interpret natural-language and source-code information. Earlier automated ODC classification systems based on SVM and Naive Bayes achieved accuracies of approximately 77--82\% and were limited by sparse textual representations \citep{huang2011autoodc, thung2012automatic}. In contrast, LLMs can integrate multiple forms of software-engineering evidence, including code structure, test behavior, and error messages, within a single inference process. Moreover, their ability to perform classification without task-specific labeled training data makes them particularly suitable for applying semantic defect taxonomies such as ODC to large-scale bug datasets.

Existing defect classification approaches primarily rely on bug reports as the source of evidence for classification. In contrast, our approach integrates multiple software artifacts, including bug reports, source code, failing tests, exceptions, and stack traces, to classify defect types. The use of diverse evidence motivates the need for a well-established and systematically defined taxonomy to constrain and guide LLM-based classification. We therefore adopt Orthogonal Defect Classification (ODC) as the underlying taxonomy for organizing and interpreting defect types.


\subsection{Research Gap}
Despite the growing adoption of LLMs in software engineering, their application to semantic defect classification using the ODC taxonomy remains largely unexplored. Existing defect classification approaches have primarily relied on individual sources of evidence, such as issue reports, code-level features, or repair patches, rather than integrating diverse software artifacts, including bug reports, source code, failing tests, exceptions, and stack traces. Furthermore, there is limited understanding of how effectively LLMs can classify defects into well-defined ODC categories and how different prompting strategies influence classification performance. The lack of ODC defect-type annotations in widely used \textit{Defects4J} benchmark also presents a significant challenge for systematic evaluation. These gaps motivate our study to investigate LLM-based ODC defect classification using multiple software artifacts and to systematically compare different prompting strategies for improving the accuracy and reliability of defect-type classification.

%% =============================================================================
